\section{Product Vacua with Boundary States}
\label{sec:examples_pvbs}

The one-species PVBS chain has a product vacuum and a geometric boundary
particle~\cite[Section~II, equations~(1)--(8)]{BachmannNachtergaele2012PVBS}.
We use left-to-right matrix products and write $q=\mu e^{i\theta}$ for the
hopping parameter and phase. The review's real parameter is $q=1+\lambda$
\cite[Section~IV.C.4]{Cirac2021Matrix}. Its printed right-localization threshold
and unconditional periodic-uniqueness claim require the corrections explained
in~\cite{gap:cpgsv21_pvbs_threshold}. The following finite-chain results allow
all complex $q$; $q=0$ is an algebraic extension of the original positive-hopping
family. No spectral-gap assertion is included.

\begin{definition}[One-species PVBS tensor and boundary vectors]
    \label{def:pvbs_tensor}
    \lean{MPSTensor.pvbsTensor}
    \lean{MPSTensor.pvbsVacuum}
    \lean{MPSTensor.pvbsEdge}
    \leanok
    Set $A^0=\operatorname{diag}(q,1)$ and $A^1=|0\rangle\langle1|$.
    Write $\Omega_N=|0\cdots0\rangle$ and
    $\Psi_N(q)=\sum_{k=0}^{N-1}q^k|0\cdots1_k\cdots0\rangle$.
\end{definition}

\begin{lemma}[PVBS amplitudes and boundary contraction]
    \label{lem:pvbs_amplitudes}
    \lean{MPSTensor.pvbsVacuum_zero}
    \lean{MPSTensor.pvbsVacuum_nil}
    \lean{MPSTensor.pvbsVacuum_cons}
    \lean{MPSTensor.pvbsEdge_cons_zero}
    \lean{MPSTensor.pvbsEdge_cons_one}
    \lean{MPSTensor.pvbsEdge_zero}
    \lean{MPSTensor.pvbsEdge_first}
    \lean{MPSTensor.pvbsVacuum_first}
    \lean{MPSTensor.evalWord_pvbsTensor}
    \lean{MPSTensor.pvbsTensor_openState}
    \lean{MPSTensor.pvbsEdge_excitedAt}
    \lean{MPSTensor.pvbsVacuum_excitedAt}
    \leanok
    \uses{def:pvbs_tensor}
    For every physical word $\sigma$,
    \begin{align}
        A^{\sigma_1}\cdots A^{\sigma_N}
         & =\begin{pmatrix}
                q^N\Omega_N(\sigma) & \Psi_N(q)(\sigma) \\
                0                   & \Omega_N(\sigma)
            \end{pmatrix}.
        \label{eq:pvbs_word}
    \end{align}
    The standard boundary contraction $\langle0|A^{\sigma_1}\cdots
        A^{\sigma_N}|1\rangle$ is $\Psi_N(q)(\sigma)$.
\end{lemma}
\begin{proof}
    \leanok
    Multiplication by $A^0$ contributes $q$ before the unique particle;
    multiplication by $A^1$ creates it. Two particles give the zero matrix.
    Induction on the word length gives (\ref{eq:pvbs_word}).
\end{proof}

\begin{theorem}[Open PVBS matrix-product space]
    \label{thm:pvbs_open_space}
    \lean{MPSTensor.groundSpaceMap_pvbsTensor}
    \lean{MPSTensor.mem_groundSpace_pvbsTensor_iff}
    \lean{MPSTensor.groundSpace_pvbsTensor}
    \lean{MPSTensor.linearIndependent_pvbs_states}
    \lean{MPSTensor.finrank_groundSpace_pvbsTensor}
    \leanok
    \uses{lem:pvbs_amplitudes}
    The full boundary matrix-product space is
    $\mathcal G_N=\operatorname{span}\{\Omega_N,\Psi_N(q)\}$.
    It has dimension two for every $N\ge1$.
\end{theorem}
\begin{proof}
    \leanok
    Taking the trace against a boundary matrix $X$ gives
    $(q^NX_{00}+X_{11})\Omega_N+X_{10}\Psi_N(q)$.
    The entries $X_{11}$ and $X_{10}$ realize arbitrary coefficients.
    Evaluation on the vacuum and on a particle at the first site proves
    linear independence.
\end{proof}

\begin{theorem}[Open PVBS parent-Hamiltonian ground space]
    \label{thm:pvbs_open_parent}
    \lean{MPSTensor.pvbs_pair_constraints}
    \lean{MPSTensor.mem_groundSpace_pvbs_two_iff}
    \lean{MPSTensor.groundSpace_pvbs_intersection}
    \lean{MPSTensor.contiguous_mem_groundSpace_pvbs}
    \lean{MPSTensor.ker_openParentHamiltonianES_pvbs}
    \lean{MPSTensor.finrank_ker_openParentHamiltonianES_pvbs}
    \leanok
    \uses{thm:pvbs_open_space}
    The open nearest-neighbour parent Hamiltonian has kernel $\mathcal G_N$
    and ground-state dimension two for every $N\ge2$.
\end{theorem}
\begin{proof}
    \leanok
    A two-site vector satisfies the local constraint precisely when
    $a_{11}=0$ and $a_{01}=q a_{10}$.
    The left and right restrictions of a putative longer ground vector
    lie in the shorter vacuum-particle spaces. Their common two-site
    constraints eliminate the two-particle coefficient and relate the
    remaining particle coefficients by $q$. This proves the intersection
    property from length two onward. Iterating it identifies all nonwrapping
    local constraints with $\mathcal G_N$. The existing positive-parent
    kernel theorem then identifies the actual Hamiltonian kernel.
\end{proof}

\begin{theorem}[Periodic PVBS vacuum uniqueness]
    \label{thm:pvbs_periodic_parent}
    \lean{MPSTensor.chainGroundSpace_pvbs_le_groundSpace}
    \lean{MPSTensor.pvbsVacuum_mem_chainGroundSpace}
    \lean{MPSTensor.cyclicTranslateCfg_excitedAt}
    \lean{MPSTensor.pvbs_periodic_particle_coefficient}
    \lean{MPSTensor.chainGroundSpace_pvbs_eq_vacuum}
    \lean{MPSTensor.ker_parentHamiltonian_pvbs}
    \lean{MPSTensor.finrank_ker_parentHamiltonian_pvbs}
    \lean{MPSTensor.pvbs_pow_ne_one_of_norm_ne_one}
    \lean{MPSTensor.ker_parentHamiltonian_pvbs_of_norm_ne_one}
    \leanok
    \uses{thm:pvbs_open_parent}
    If $N\ge2$ and $q^N\ne1$, the periodic parent Hamiltonian has kernel
    $\operatorname{span}\{\Omega_N\}$ and ground-state dimension one.
    In particular, this conclusion holds at every such length if $|q|\ne1$.
\end{theorem}
\begin{proof}
    \leanok
    Every periodic ground vector satisfies the nonwrapping constraints,
    hence is $a\Omega_N+b\Psi_N(q)$. A cyclic translate is also a ground
    vector. Comparing its particle amplitudes at the first two sites yields
    $b=bq^N$, forcing $b=0$. The vacuum satisfies every cyclic constraint.
\end{proof}

\begin{theorem}[Critical W state and geometric localization]
    \label{thm:pvbs_boundary_regimes}
    \lean{MPSTensor.pvbsTensor_one}
    \lean{MPSTensor.pvbsEdge_one}
    \lean{MPSTensor.pvbsEdge_one_two_mem_chainGroundSpace}
    \lean{MPSTensor.pvbs_critical_periodic_counterexample}
    \lean{MPSTensor.norm_pvbsEdge_excitedAt}
    \lean{MPSTensor.pvbsEdge_left_decay}
    \lean{MPSTensor.pvbsEdge_right_relative}
    \lean{MPSTensor.pvbsEdge_right_decay}
    \leanok
    \uses{lem:pvbs_amplitudes,thm:pvbs_periodic_parent}
    At $q=1$, the tensor and boundary vector are the W-state tensor and
    unnormalized W vector. Already on the two-site ring, the W vector is a
    nonvacuum zero-energy state. For $|q|<1$, amplitudes decay as $|q|^k$
    from the left edge. For $|q|>1$, the amplitude $k$ sites from the right
    edge relative to the last-site amplitude equals $q^{-k}$ and tends to zero.
\end{theorem}
\begin{proof}
    \leanok
    Substitution gives the W tensor at $q=1$ and verifies the two periodic
    constraints. Evaluation on a one-particle configuration separates the
    W vector from the vacuum span. The localization statements follow from
    the exact geometric coefficients and convergence of powers of a scalar
    of modulus less than one.
\end{proof}
